\scanpage{198}
\mspage{54}
Or $I'_{/u^*(i)}\simeq I'_{/i}$, \deleted{\unsure{do}} donc on a la flèche

\[
 (u_!u^*(i))\unsure{j)|(j)}\simeq\pi_0\bigl({}_{j\backslash}(I'_{/i})\bigr)\longrightarrow\operatorname{Hom}_I(j,i)
 \tag*{(\unsure{4|5}.10)}
\]

Or on voit aussitôt que pour $\alpha\colon j\longrightarrow i$ dans $\operatorname{Hom}(j,i)$, la fibre de (4.10) en $u$ s'identifie à $\pi_0({}_{j\backslash}\underset{\alpha}{I'}_{/i})$, de façon plus précise, $\,{}_{j\backslash}\underset{\alpha}{I'}_{/i}$ est la sous-catégorie ouverte et fermée de $J={}_{j\backslash}(I'_{/i})$ \unsure{correspondant} à la partie de $\pi_0(J)$ \deleted{fibre} qui est au dessus de $\alpha$ dans (4.10.). \deleted{\unsure{Cel}} \inserted{Cela achève} la \deleted{\unsure{signification}} \inserted{démonstration} de 4.1.

\deleted{achève la démonstration de 4.1.}

NB Une autre façon de prouver que $c\Rightarrow A$ c'est de noter que les deux mbs du \unsure{le} morphisme d'adjonction $u_!u^*F\mathbin{\unsure{-|\to}}F$ dans $I^\wedge$ sont des foncteurs $I^\wedge\mathbin{-}I^\wedge$ qui commutent aux $\unsure{\varinjlim|\underline{\lim}}$ quelconques. Donc pour vérifier que \deleted{\unsure{des}} c'est iso, il suffit de le vérifier pour $F$ dans $I$, puisque $I$ engendre $I^\wedge$ par $\varinjlim$.

Par contre, \deleted{\unsure{cet}} \inserted{un} argument \inserted{de ce type} ne semble pas prouver que pour vérifier que l'on a un \uline{iso} $F\to u_*u^*F$, il suffise de le faire pour $F$ représentable, \deleted{\unsure{Alors}} i.e. de vérifier \unsure{$b''$)|$d''$)} pour $F$ représentable. \deleted{Il nous reste \unsure{donc}} Il nous reste à \deleted{prouver} \inserted{voir si} $c'\Rightarrow d''$.
% archive sheet number: 198
